\chapter{重积分}\label{ch:7}
\begin{tcolorbox}[title=章首知识地图,colframe=Ink]
区域与体积元 $\to$ 累次积分 $\to$ 换序 $\to$ 极、柱、球坐标 $\to$ 一般换元、对称与应用。
\end{tcolorbox}
\begin{examtip}考研主线：二重三重积分、换序、极柱球坐标及应用。常见题型：读图定界、换元化简、对称消项。竞赛拓展：Gaussian积分、一般Jacobian和积分不等式；换序前确认收敛条件。上述为学习任务分层，具体范围以所报考科目的当年大纲为准。\end{examtip}
\section{重积分定义与累次积分}\label{sec:7:1}
\subsection{概念、定理与推导}

\begin{definition}{二重、三重Riemann积分}
设$D$为有界Jordan可测集，即其边界的面积为零，且$f:D\to\R$有界。将$D$分成有限个Jordan可测小块$D_i$，各块任取$\xi_i$，作和$\sum_i f(\xi_i)|D_i|$。若当最大块直径趋零时，这些和不论如何分割和取点都趋于同一有限值，就称$f$在$D$上Riemann可积，极限记为$\iint_Df\dd A$。等价地，可在包含$D$的矩形上将$f$零延拓后定义积分。
三维用有界Jordan可测集、体积及$\sum_i f(\xi_i)|\Omega_i|$同样定义$\iiint_\Omega f\dd V$。常函数1给面积或体积；无界区域、无界函数的积分需另取反常极限。本书涉及Tonelli、几乎处处和$L^p$时采用Lebesgue积分，不能将其可积条件与本定义混同。
\end{definition}
\begin{theorem}{Fubini的常用版本，标准结论}若$f$在由分段光滑边界围成的有界闭区域上连续，则可按区域投影写成累次积分；例如$D=\{a\le x\le b,\varphi_1(x)\le y\le\varphi_2(x)\}$时，$\iint_Df=\int_a^b\dd x\int_{\varphi_1(x)}^{\varphi_2(x)}f\dd y$。\end{theorem}
\begin{proof}先沿条带分割，条带内的Riemann和趋于内层积分，再对条带宽度求和。连续性和紧致性保证误差统一趋零。反常积分换序需绝对可积或非负等额外条件。\end{proof}
重积分具有线性、区域可加性、保序性；若$m\le f\le M$，则$m|D|\le\iint_Df\le M|D|$。

\begin{method}{重积分定义与累次积分}条件：先描绘区域。步骤：投影决定外限，穿过区域的切线段决定内限；必要时分块。易错点：内限可依赖外变量，外限不得依赖内变量。\end{method}
\subsection{例题与方法比较}
\begin{example}{基础例题｜矩形}{难度：$\star$}
\label{ex:7:1}\textbf{题目} 求$\int_0^1\dd x\int_0^2(x+y)\dd y$。\par
\textbf{思路} 直接积分或分离。\par
\textbf{解答} 方法一：内层为$2x+2$，外层得3。方法二：$\iint x=2\int_0^1x\dd x=1$，$\iint y=\int_0^2y\dd y=2$，相加3。\par
\textbf{方法点评} 可分离项利用矩形区域简化。\par
\textbf{易错点或变式} 积分区域与被积式一起决定能否拆积。
\end{example}
\begin{example}{考研风格例题｜三角区域}{难度：$\star\star$}
\label{ex:7:2}\textbf{题目} $D=\{x,y\ge0,x+y\le1\}$，求$\iint_D(x+y)\dd A$。\par
\textbf{思路} 累次积分或利用对称与重心。\par
\textbf{解答} 方法一：$\int_0^1\dd x\int_0^{1-x}(x+y)\dd y=\frac{1}{3}$。方法二：对称给两项相等，$\iint_Dx=\int_0^1x(1-x)\dd x=\frac{1}{6}$，总和$\frac{1}{3}$。\par
\textbf{方法点评} 交换$x,y$的对称性可将计算减半。\par
\textbf{易错点或变式} 三角区域不允许简单拆成两个独立的一维积分乘积。
\end{example}
\begin{example}{考研风格例题｜面积}{难度：$\star\star$}
\label{ex:7:3}\textbf{题目} 求抛物线$y=x^2$与直线$y=x$之间面积。\par
\textbf{思路} 按$x$或按$y$切片。\par
\textbf{解答} 方法一：$0\le x\le1$，面积$\int_0^1(x-x^2)\dd x=\frac{1}{6}$。方法二：$0\le y\le1$，$y\le x\le\sqrt y$，面积$\int_0^1(\sqrt y-y)\dd y=\frac{1}{6}$。\par
\textbf{方法点评} 换序能从另一方向检验边界。\par
\textbf{易错点或变式} 解$x^2\le y$时本题位于第一象限。
\end{example}
\section{换序与极坐标}\label{sec:7:2}

\begin{center}\begin{tikzpicture}[scale=2.4,>=stealth]
\fill[Teal!12](0,0)--(0,1)--(1,1)--cycle;
\draw[->,Ink](0,0)--(1.25,0)node[right]{$x$};\draw[->,Ink](0,0)--(0,1.2)node[above]{$y$};
\draw[Teal,thick](0,0)--(1,1)--(0,1);\draw[Gold,thick](.4,.4)--(.4,1);\draw[Ink,thick](0,.7)--(.7,.7);
\node[right] at(.85,.5){$x=y$};\node[left]at(0,1){$1$};\node[below]at(1,0){$1$};
\end{tikzpicture}\par\small 同一区域：竖切为$x\le y\le1$，横切为$0\le x\le y$。
\end{center}

\subsection{概念、定理与推导}

\begin{theorem}{极坐标换元}令$x=r\cos\theta,y=r\sin\theta$，则$\dd A=r\dd r\dd\theta$，$r\ge0$。对有界可积函数及适当分块区域，重积分可按此换元。\end{theorem}
\begin{proof}雅可比$\det\left(\frac{\partial(x,y)}{\partial(r,\theta)}\right)=r$，也可由小扇形面积近似$r\Delta r\Delta\theta$得到。原点退化是面积为零的集合，不影响通常积分。\end{proof}
换序不改变区域，只改变描述。若原顺序出现$\e^{y^2}$等无初等原函数，尝试先对另一变量积分。圆、扇形、圆环与$x^2+y^2$结构通常适合极坐标。

\begin{method}{换序与极坐标}条件：角域须覆盖区域且内部不重复。步骤：先定角，再由射线交点定半径。易错点：漏掉雅可比$r$或把偏心圆误写成常半径。\end{method}
\subsection{例题与方法比较}
\begin{example}{考研风格例题｜换序消障碍}{难度：$\star\star$}
\label{ex:7:4}\textbf{题目} 求$\int_0^1\dd x\int_x^1\e^{y^2}\dd y$。\par
\textbf{思路} 区域为$0\le x\le y\le1$。\par
\textbf{解答} 方法一：换序得$\int_0^1\dd y\int_0^y\e^{y^2}\dd x=\int_0^1y\e^{y^2}\dd y=\frac{\e-1}{2}$。方法二：令$F(x)=\int_x^1\e^{y^2}\dd y$，分部积分$\int_0^1F\dd x=[xF]_0^1-\int_0^1xF\prime\dd x$也得同式。\par
\textbf{方法点评} 换序最直观，分部法从变限积分得到同一权重。\par
\textbf{易错点或变式} “无初等原函数”不代表定积分不可计算。
\end{example}
\begin{example}{考研风格例题｜圆盘二次矩}{难度：$\star\star$}
\label{ex:7:5}\textbf{题目} 单位圆盘$D$上求$\iint_D(x^2+y^2)\dd A$。\par
\textbf{思路} 径向函数与圆域匹配。\par
\textbf{解答} 方法一：极坐标$2\pi\int_0^1r^3\dd r=\frac{\pi}{2}$。方法二：薄圆环面积$2\pi r\dd r$，乘$r^2$积分也得$\frac{\pi}{2}$；直接直角坐标可写$\int_{-1}^1\int_{-\sqrt{1-x^2}}^{\sqrt{1-x^2}}(x^2+y^2)\dd y\dd x$验证。\par
\textbf{方法点评} 极坐标是圆环分层的坐标表达。\par
\textbf{易错点或变式} $r^2$来自被积式，另一个$r$来自面积元。
\end{example}
\begin{example}{考研风格例题｜偏心圆}{难度：$\star\star$}
\label{ex:7:6}\textbf{题目} 求区域$x^2+y^2\le2x$的面积。\par
\textbf{思路} 用平移圆或极坐标。\par
\textbf{解答} 配方$(x-1)^2+y^2\le1$，面积$\pi$。极坐标为$-\frac{\pi}{2}\le\theta\le\frac{\pi}{2},0\le r\le2\cos\theta$，面积$\frac12\int_{-\frac{\pi}{2}}^{\frac{\pi}{2}}4\cos^2\theta\dd\theta=\pi$。\par
\textbf{方法点评} 配方快速识图，极坐标适合附加径向权函数。\par
\textbf{易错点或变式} 半径上界不能取负，故角度有约束。
\end{example}
\section{三重积分与柱、球坐标}\label{sec:7:3}
\subsection{概念、定理与推导}

\begin{definition}{柱面与球面坐标}柱坐标$x=r\cos\theta,y=r\sin\theta,z=z$，体积元$r\dd r\dd\theta\dd z$。球坐标$x=\rho\sin\phi\cos\theta,y=\rho\sin\phi\sin\theta,z=\rho\cos\phi$，体积元$\rho^2\sin\phi\dd\rho\dd\phi\dd\theta$，$0\le\phi\le\pi$。\end{definition}
\begin{theorem}{分层与投影}若区域在每个$(x,y)$上有$z_1\le z\le z_2$，则先对$z$积分再投影到$xy$平面；若固定$z$截面更简单，则先对截面作二重积分再对$z$积分。\end{theorem}
\begin{proof}两种分层方式均可由积分和重新分组得到。球坐标下三个相互垂直方向的小边长分别为\[\dd\rho,\qquad\rho\dd\phi,\qquad\rho\sin\phi\dd\theta.\]三者相乘得到体积元。\end{proof}

\begin{method}{三重积分与柱、球坐标}条件：坐标约定明确。步骤：圆柱对称用柱坐标，球或锥与球相交用球坐标；画经线剖面定范围。易错点：将极角与纬度混用导致正弦余弦错误。\end{method}
\subsection{例题与方法比较}
\begin{example}{基础例题｜球体积}{难度：$\star$}
\label{ex:7:7}\textbf{题目} 求半径$R$的球体积。\par
\textbf{思路} 球坐标或截面。\par
\textbf{解答} 球坐标积分$\int_0^{2\pi}\dd\theta\int_0^\pi\sin\phi\dd\phi\int_0^R\rho^2\dd\rho=\frac{4\pi R^3}{3}$。截面法$\int_{-R}^R\pi(R^2-z^2)\dd z$给同值。\par
\textbf{方法点评} 两种分层方向检验体积元。\par
\textbf{易错点或变式} 球面面积与球体积的幂次不同。
\end{example}
\begin{example}{考研风格例题｜锥体}{难度：$\star\star$}
\label{ex:7:8}\textbf{题目} $0\le z\le1,x^2+y^2\le z^2$，求体积。\par
\textbf{思路} 固定$z$为半径$z$的圆盘。\par
\textbf{解答} $V=\int_0^1\pi z^2\dd z=\frac{\pi}{3}$。柱坐标也可取$0\le r\le1,r\le z\le1$，得$2\pi\int_0^1r(1-r)\dd r=\frac{\pi}{3}$。\par
\textbf{方法点评} 换序对应从截圆改为竖条柱。\par
\textbf{易错点或变式} 锥面是$r=z$，不是$r=z^2$。
\end{example}
\begin{example}{考研风格例题｜球内二次矩}{难度：$\star\star$}
\label{ex:7:9}\textbf{题目} 求单位球内$\iiint z^2\dd V$。\par
\textbf{思路} 旋转对称或球坐标。\par
\textbf{解答} 对称使三坐标平方积分相等，总和为$\iiint\rho^2\dd V=\frac{4\pi}{5}$，故每项$\frac{4\pi}{15}$。直接球坐标的角积分为$\int_0^{2\pi}\dd\theta\int_0^\pi\cos^2\phi\sin\phi\dd\phi=\frac{4\pi}{3}$，乘径向积分$\frac{1}{5}$同得。\par
\textbf{方法点评} 对称先降复杂度。\par
\textbf{易错点或变式} 被积函数偶对称意味着相等，不意味着积分为零。
\end{example}
\section{一般换元、Jacobian与应用}\label{sec:7:4}

\begin{method}{从区域、密度到物理量}
以常密度求质心时先用对称性确定为零或相等的坐标，再算一个必要的一阶矩。平面薄片的质量为$\iint_D\sigma\dd A$，绕原点的极惯量为$\iint_D(x^2+y^2)\sigma\dd A$。空间体绕$x,y,z$轴分别用到$y^2+z^2,x^2+z^2,x^2+y^2$。总质量必须非零，惯量使用到轴的距离平方，不能使用沿轴坐标平方。
\end{method}

\subsection{概念、定理与推导}

\begin{theorem}{多元换元，标准结论}若$T$在区域内部一一对应、连续可微且Jacobian非零，边界可忽略，则
\begin{equation}\iint_Df(x,y)\dd x\dd y=\iint_G f(T(u,v))\left|\dfrac{\partial(x,y)}{\partial(u,v)}\right|\dd u\dd v.\end{equation}
三维公式同理。线性变换的面积或体积倍率为行列式绝对值。\end{theorem}
\begin{proof}在小块上以导数矩阵线性近似，行列式给局部体积倍率；累加并控制误差得到换元公式。\end{proof}
质量$M=\iiint\rho\dd V$，质心$\bar x=M^{-1}\iiint x\rho\dd V$，绕$z$轴转动惯量$I_z=\iiint(x^2+y^2)\rho\dd V$。密度与球坐标径向变量应按上下文区分。

\begin{method}{一般换元、Jacobian与应用}条件：须确认变换可逆且覆盖原区域。步骤：由边界线性组合选$u,v$，反解$x,y$，计算反向Jacobian。易错点：使用正向行列式却没有取倒数。\end{method}
\subsection{例题与方法比较}
\begin{example}{考研风格例题｜椭圆}{难度：$\star\star$}
\label{ex:7:10}\textbf{题目} 求椭圆$\frac{x^2}{a^2}+\frac{y^2}{b^2}\le1$面积，$a,b>0$。\par
\textbf{思路} 伸缩到单位圆。\par
\textbf{解答} $x=au,y=bv$，Jacobian为$ab$，面积$ab\pi$。也可令$x=ar\cos\theta,y=br\sin\theta$，面积元$abr\dd r\dd\theta$，积分同得。\par
\textbf{方法点评} 线性伸缩与极坐标可分两步或合并。\par
\textbf{易错点或变式} 不能把椭圆周长公式误当成面积公式。
\end{example}
\begin{example}{考研风格例题｜斜域换元}{难度：$\star\star$}
\label{ex:7:11}\textbf{题目} $D$由$0\le x+y\le1,0\le x-y\le1$定义，求面积。\par
\textbf{思路} 两个线性组合把平行四边形变成正方形。\par
\textbf{解答} 令$u=x+y,v=x-y$，则$x=\frac{u+v}{2},y=\frac{u-v}{2}$，反向Jacobian绝对值$\frac{1}{2}$。因此面积$\int_0^1\int_0^1\frac12\dd u\dd v=\frac{1}{2}$。\par
\textbf{方法点评} 区域边界常比被积式更能提示换元。\par
\textbf{易错点或变式} 正向Jacobian为$-2$，不能直接乘2。
\end{example}
\begin{example}{考研风格例题｜质心}{难度：$\star\star$}
\label{ex:7:12}\textbf{题目} 均匀三角形$D=\{x,y\ge0,x+y\le1\}$的质心。\par
\textbf{思路} 用一阶矩或顶点平均。\par
\textbf{解答} 面积$\frac{1}{2}$，$\iint_Dx\dd A=\frac{1}{6}$，故$\bar x=\frac{1}{3}$；对称给$\bar y=\frac{1}{3}$。三顶点平均也得$(\frac{1}{3},\frac{1}{3})$。\par
\textbf{方法点评} 积分法可推广非均匀密度，几何公式限均匀三角形。\par
\textbf{易错点或变式} 质心坐标需除以总质量。
\end{example}
\section{对称、Gaussian积分与综合}\label{sec:7:5}
\subsection{概念、定理与推导}

\begin{theorem}{重积分对称性}若区域关于$x\mapsto-x$对称，被积式关于$x$为奇，则积分为零；若为偶，可在半区域积分乘2。置换$x,y$保持区域和体积元时，两置换后的积分相等。\end{theorem}
\begin{proof}在换元公式中使用反射或坐标置换，Jacobian绝对值为1，奇函数积分等于其相反数。\end{proof}
\begin{extension}非负积分允许在适当极限下交换积分次序；绝对可积也允许换序。Gaussian积分通过平方转为二维径向积分，是把一维困难问题提升维数后简化的典型。\end{extension}

\begin{method}{对称、Gaussian积分与综合}条件：区域和被积式都检查对称性；无穷区域先判非负或绝对可积。步骤：先消奇项，再统一同类项，最后选坐标。易错点：仅函数奇而区域不对称不能消去。\end{method}
\subsection{例题与方法比较}
\begin{example}{竞赛风格例题｜Gaussian积分}{难度：$\star\star\star$}
\label{ex:7:13}\textbf{题目} 求$I=\int_{-\infty}^{\infty}\e^{-x^2}\dd x$。\par
\textbf{思路} 平方后极坐标。\par
\textbf{解答} 积分正且有限（尾部可与$\e^{-|x|}$比较）。由非负积分换序，$I^2=\iint_{\R^2}\e^{-x^2-y^2}\dd A=2\pi\int_0^\infty r\e^{-r^2}\dd r=\pi$，故$I=\sqrt\pi$。\par
\textbf{方法点评} 升维使指数之和成为半径平方。\par
\textbf{易错点或变式} 最后取正根，因为$I>0$。
\end{example}
\begin{example}{竞赛风格例题｜奇偶消项}{难度：$\star\star\star$}
\label{ex:7:14}\textbf{题目} 单位球内求$\iiint(x^3+xy+z^2)\dd V$。\par
\textbf{思路} 先用反射对称。\par
\textbf{解答} $x^3$对$x$奇，$xy$也对$x$奇，积分均零；剩$z^2$积分为$\frac{4\pi}{15}$。无需处理三项的完整球坐标表达式。\par
\textbf{方法点评} 先对称再坐标可避免大量角积分。\par
\textbf{易错点或变式} 球必须以原点为中心才有这里使用的反射对称。
\end{example}
\begin{example}{竞赛风格例题｜球锥交域}{难度：$\star\star\star$}
\label{ex:7:15}\textbf{题目} 求单位球内且$z\ge\sqrt{x^2+y^2}$的体积。\par
\textbf{思路} 球坐标中锥条件为$0\le\phi\le\frac{\pi}{4}$。\par
\textbf{解答} $V=\int_0^{2\pi}\dd\theta\int_0^{\frac{\pi}{4}}\sin\phi\dd\phi\int_0^1\rho^2\dd\rho=\frac{2\pi}{3}(1-\frac{1}{\sqrt2})$。\par
\textbf{方法点评} 球面与锥面的交域在球坐标中边界全为常数。\par
\textbf{易错点或变式} 锥条件同时要求$z\ge0$，不能再加下半球。
\end{example}
\begin{warning}重积分换元中的Jacobian取绝对值；坐标覆盖次数必须正确；有符号函数在无穷区域换序前需要额外收敛条件。\end{warning}


\section{重积分的换元证明、广义积分与推广}\label{mat:multiple}
\begin{examtip}进一步处理Jacobian、积分范围、反常积分与高维球。\end{examtip}
\subsection{Jacobian与局部面积倍率}\label{mat:multiple:jacproof}
\begin{method}{Jacobian与局部面积倍率}\textbf{适用条件} $C^1$局部可逆变换。\par\textbf{实施与依据} 小矩形的两条边由导数矩阵线性映射，行列式给有向面积倍率。\end{method}
\begin{example}{专题例题｜Jacobian与局部面积倍率}{难度：$\star\star$}
\textbf{题目} 用$x=au\cos v,y=bu\sin v$计算椭圆二次矩$\iint x^2\dd A$。\par\textbf{思路与解答} Jacobian为$abu$，$0\le u\le1,0\le v\le2\pi$；积分为$a^3b\int_0^1u^3\dd u\int_0^{2\pi}\cos^2v\dd v=\frac{\pi a^3b}{4}$。\par\textbf{方法点评与边界} 并非把圆盘结果中的半径简单替换为一个半轴。
\end{example}
\subsection{Fubini与Tonelli的区别}\label{mat:multiple:tonelli}
\begin{method}{Fubini与Tonelli的区别}\textbf{适用条件} 无穷区域或有符号函数。\par\textbf{实施与依据} 非负函数可用Tonelli，有限绝对积分可用Fubini；否则不能保证换序。\end{method}
\begin{example}{专题例题｜Fubini与Tonelli的区别}{难度：$\star\star$}
\textbf{题目} 在第一象限计算$\iint\e^{-x-y}\dd x\dd y$。\par\textbf{思路与解答} 函数非负，两种顺序均可；积分分离成$\int_0^\infty\e^{-x}\dd x\int_0^\infty\e^{-y}\dd y=1$。此处也绝对可积。\par\textbf{方法点评与边界} 局部可积不等于全域绝对可积。
\end{example}
\subsection{高维球与Gamma}\label{mat:multiple:nball}
\begin{method}{高维球与Gamma}\textbf{适用条件} $d$维Euclid体积。\par\textbf{实施与依据} Gaussian积分的乘积与径向积分比较。\end{method}
\begin{example}{专题例题｜高维球与Gamma}{难度：$\star\star$}
\textbf{题目} 推导单位$d$维球体积$V_d$。\par\textbf{思路与解答} $\int_{\R^d}\e^{-|x|^2}\dd x=\pi^{\frac{d}{2}}$；径向积分为$dV_d\int_0^\infty r^{d-1}\e^{-r^2}\dd r=V_d\Gamma(\frac{d}{2}+1)$，故$V_d=\frac{\pi^{\frac{d}{2}}}{\Gamma(\frac{d}{2}+1)}$。\par\textbf{方法点评与边界} 球面面积为$dV_d$，体积与面积不能混淆。
\end{example}
\subsection{重积分的奇点判别}\label{mat:multiple:singular}
\begin{method}{重积分的奇点判别}\textbf{适用条件} 孤立径向奇性$|x|^{-p}$。\par\textbf{实施与依据} 把体积元幂次合并后比较一维反常幂积分。\end{method}
\begin{example}{专题例题｜重积分的奇点判别}{难度：$\star\star$}
\textbf{题目} 判断三维单位球内$\iiint|\bm x|^{-p}\dd V$。\par\textbf{思路与解答} 化成$4\pi\int_0^1r^{2-p}\dd r$，当且仅当$p<3$收敛，值$\frac{4\pi}{3-p}$。\par\textbf{方法点评与边界} 临界指数由维数决定，不是统一的$p<1$。
\end{example}




\section{重积分的换序、变换、应用与含参积分}\label{deep:multiple_source_deep}

\begin{examtip}补充积分定义、换序、Jacobian、几何应用、近似计算与参数极限的完整算例。\end{examtip}

\subsection{区域描述、换序与对称性}\label{deep:multiple_source_deep:region}

\begin{method}{区域描述、换序与对称性}\textbf{适用条件} 先画投影或确定边界交点，再写积分限。\par\textbf{步骤与依据} 内层积分困难时交换次序；被积函数只依赖半径时改极坐标；非矩形区域不能随意拆乘积。\par\textbf{易错点} 换序会改变积分上下限，不能只交换微分符号。\end{method}

\begin{example}{考研风格｜无初等内层积分的换序}{难度：$\star\star\star$}\label{deep:ex:multiple_source_deep:order}\textbf{题目} 求$\int_0^1\int_x^1\e^{y^2}\dd y\dd x$。\par\textbf{思路} 区域$0\le x\le y\le1$，先对$x$积分。\par\textbf{详细解答} 换序为$\int_0^1\int_0^y\e^{y^2}\dd x\dd y=\int_0^1y\e^{y^2}\dd y=\frac{\e-1}{2}$。另一法将原积分写成$F(x)=\int_x^1\e^{y^2}\dd y$的一维分部积分，$\int_0^1F(x)\dd x=[xF(x)]_0^1-\int_0^1xF'(x)\dd x$，同得结果。\par\textbf{方法点评} 换序与分部在三角区域中常是同一“累积长度”机制。\par\textbf{变式与易错点} 不必为$\int\e^{y^2}\dd y$引入特殊函数。\end{example}

\begin{example}{考研风格｜平行四边形的线性换元}{难度：$\star\star\star$}\label{deep:ex:multiple_source_deep:linearjac}\textbf{题目} $D$由$0\le x+y\le2,0\le x-y\le1$围成，求$\iint_D(x+y)^2\dd A$。\par\textbf{思路} $u=x+y,v=x-y$把区域拉直。\par\textbf{详细解答} $x=\frac{u+v}{2},y=\frac{u-v}{2}$，Jacobian绝对值$\frac{1}{2}$。所以积分$=\frac12\int_0^2u^2\dd u\int_0^1\dd v=\frac{4}{3}$。直接法需按顶点分段写上下限，结果一致但较长。变换后的面积为$\frac{2\cdot1\cdot1}{2}=1$，也可作为Jacobian核对。\par\textbf{方法点评} 边界的线性组合是最强的换元提示。\par\textbf{变式与易错点} $\frac{\partial(u,v)}{\partial(x,y)}=-2$，积分中要用逆Jacobian的绝对值$\frac{1}{2}$。\end{example}

\subsection{三重积分、椭球、球坐标与高维推广}\label{deep:multiple_source_deep:coordinates}

\begin{method}{三重积分、椭球、球坐标与高维推广}\textbf{适用条件} 根据区域对称性选择柱面、球面或线性缩放。\par\textbf{步骤与依据} 先写体元；半径、极角和方位角范围分别确定；对奇函数先使用对称性。\par\textbf{易错点} 球体元的$\rho^2\sin\phi$不可漏。\end{method}

\begin{example}{考研风格｜椭球体积与二阶矩}{难度：$\star\star\star$}\label{deep:ex:multiple_source_deep:ellipsoid}\textbf{题目} 椭球$\frac{x^2}{a^2}+\frac{y^2}{b^2}+\frac{z^2}{c^2}\le1$，求体积及$\iiint x^2\dd V$，$a,b,c>0$。\par\textbf{思路} 缩放到单位球，再球坐标或用对称性。\par\textbf{详细解答} $x=au,y=bv,z=cw$，体元$abc\dd u\dd v\dd w$。体积$\frac{4\pi abc}{3}$。对称性给$\int_Bu^2=\frac13\int_B(u^2+v^2+w^2)=\frac13 4\pi\int_0^1r^4\dd r=\frac{4\pi}{15}$，所以二阶矩$\frac{4\pi a^3bc}{15}$。直接球坐标角积分也给同值。\par\textbf{方法点评} 对称性减少角积分，缩放负责所有长度因子。\par\textbf{变式与易错点} 质量二阶矩与转动惯量不同，绕$x$轴惯量积分的是$y^2+z^2$。\end{example}

\begin{example}{竞赛拓展｜高维球体积与Gaussian}{难度：$\star\star\star$}\label{deep:ex:multiple_source_deep:nball}\textbf{题目} 推导$n$维单位球体积$V_n$。\par\textbf{思路} 分别用直角坐标乘积和径向积分计算$\int_{\R^n}\e^{-|x|^2}\dd x$。\par\textbf{详细解答} 直角坐标给$\pi^{\frac{n}{2}}$。径向层壳给$nV_n\int_0^\infty r^{n-1}\e^{-r^2}\dd r=(\frac{nV_n}{2})\Gamma(\frac{n}{2})$。所以$V_n=\frac{\pi^{\frac{n}{2}}}{\Gamma(\frac{n}{2}+1)}$。也可切片得$V_n=V_{n-1}\int_{-1}^1(1-t^2)^{\frac{n-1}{2}}\dd t$，用Beta函数递推核对。\par\textbf{方法点评} 这是$n$重积分与特殊函数的自然连接。\par\textbf{变式与易错点} 单位球表面积为$nV_n$，两者差一个$n$。\end{example}

\subsection{几何、物理与积分不等式}\label{deep:multiple_source_deep:applications}

\begin{method}{几何、物理与积分不等式}\textbf{适用条件} 先说明积分元代表面积、质量、功还是弧长。\par\textbf{步骤与依据} 从小块累积建立公式，再用对称性确定质心或消去奇矩。\par\textbf{易错点} 密度不是常数时不能以几何面积替代质量。\end{method}

\begin{example}{考研风格｜半圆薄片的质心}{难度：$\star\star\star$}\label{deep:ex:multiple_source_deep:centroid}\textbf{题目} 均匀上半圆盘$x^2+y^2\le R^2,y\ge0$的质心在哪里？\par\textbf{思路} 对称性给$\bar x=0$，极坐标算$y$的一阶矩。\par\textbf{详细解答} 面积$A=\frac{\pi R^2}{2}$。$\iint_Dy\dd A=\int_0^\pi\int_0^R(r\sin\theta)r\dd r\dd\theta=\frac{2R^3}{3}$，所以$\bar y=\frac{4R}{3\pi}$。按水平切片也可算$\int_0^R2y\sqrt{R^2-y^2}\dd y$，同得一阶矩。\par\textbf{方法点评} 两坐标系给相同结果；先对称后计算可减少工作。\par\textbf{变式与易错点} 半圆弧的质心为$\frac{2R}{\pi}$，与半圆薄片不同。\end{example}

\begin{example}{竞赛风格｜积分Cauchy与等号条件}{难度：$\star\star\star$}\label{deep:ex:multiple_source_deep:holder}\textbf{题目} 设$a<b$且$f,g$为区间上实平方可积函数，证明$(\int_a^bf g)^2\le\int_a^bf^2\int_a^bg^2$，并求等号条件。\par\textbf{思路} 非负二次多项式判别式，或最小平方距离。\par\textbf{详细解答} 令$A=\int f^2,B=\int fg,C=\int g^2$。若$C>0$，对所有$t$有$\int(f-tg)^2=A-2Bt+Ct^2\ge0$，其最小值$A-\frac{B^2}{C}\ge0$，即结论。等号当且仅当$f=(\frac{B}{C})g$几乎处处。若$C=0$，则$g=0$几乎处处，结论自动成立。\par\textbf{方法点评} 证明同时给精确等号条件，可用于最值和积分估计。\par\textbf{变式与易错点} “线性相关”是几乎处处意义，连续函数时可提升为处处。\end{example}

\subsection{积分近似与参数收敛}\label{deep:multiple_source_deep:numerics}

\begin{method}{积分近似与参数收敛}\textbf{适用条件} 数值公式需要相应阶导数有界；参数换序需要统一控制。\par\textbf{步骤与依据} 先选择梯形、Simpson或Gauss型求积，再给误差阶；函数列换积分先证一致或控制收敛。\par\textbf{易错点} 仅有点态收敛不能保证积分极限相等。\end{method}

\begin{example}{竞赛拓展｜求积公式的精度检查}{难度：$\star\star\star$}\label{deep:ex:multiple_source_deep:simpson}\textbf{题目} 用单段Simpson公式近似$\int_0^1\e^x\dd x$并给误差界。\par\textbf{思路} 公式对三次以下多项式精确，四阶导数控制余项。\par\textbf{详细解答} $S=\frac{f(0)+4f(\frac{1}{2})+f(1)}{6}=\frac{1+4\sqrt\e+\e}{6}$。标准误差为$\int_0^1f-S=-\frac{f^{(4)}(\xi)}{2880}$，故$|\int_0^1\e^x-S|\le\frac{\e}{2880}$。真实值$\e-1$，可直接比较。作为核对，对$f=x^4$有积分$\frac{1}{5}$、Simpson值$\frac{5}{24}$，差$-\frac{1}{120}=-\frac{24}{2880}$，符号及常数吻合。\par\textbf{方法点评} 误差式属于标准求积结论，四阶余项或Peano核可证明；本例给独立多项式核对。\par\textbf{变式与易错点} 复合Simpson要求等分段数为偶数。\end{example}

\begin{example}{竞赛风格｜积分与极限失效的尖峰}{难度：$\star\star\star$}\label{deep:ex:multiple_source_deep:badlimit}\textbf{题目} 取$f_n(x)=n\mathbf1_{(0,\frac{1}{n})}(x)$，$x\in[0,1]$。比较点态极限与积分极限。\par\textbf{思路} 高度增大、宽度减小但面积不变。\par\textbf{详细解答} 每个固定$x>0$最终在支撑外，且$f_n(0)=0$，所以逐点趋0。但$\int_0^1f_n=1$，不趋0。没有共同可积控制函数：若有$g\ge f_n$，则控制收敛将推出积分趋0，矛盾。若改成$g_n=\mathbf1_{(0,\frac{1}{n})}$，受1控制，积分$\frac{1}{n}\to0$。\par\textbf{方法点评} 用一个模型同时解释一致收敛、控制收敛与质量集中。\par\textbf{变式与易错点} “区间有限”本身不足以交换点态极限与积分。\end{example}

\section{重积分的轮换配对与齐次边界}\label{ic:delta_multiple}

\begin{examtip}本节在奇偶性、换序、柱面与球面坐标的基础上，讨论函数配对、区域构型和计算方法比较。\end{examtip}

\subsection{轮换后配成径向函数}\label{ic:delta_multiple:permutationpair}\begin{method}{轮换后配成径向函数}\textbf{识别信号} 区域在$x,y$互换下不变，被积函数互换后可用三角恒等式合并。\par \textbf{操作步骤} 先写$I=\iint f(y,x)$，两式相加除以2，再选极坐标。\par \textbf{疑难点与条件} 区域对称与被积函数配对是两项独立条件。\end{method}

\begin{example}{竞赛风格｜非径向被积式的对称化}{难度：$\star\star\star$}\label{ic:ex:delta_multiple:sinepair}\textbf{题目} 求单位圆盘上$\iint\sin(x^2)\cos(y^2)\dd A$。\par \textbf{思路} 交换$x,y$后用正弦和角公式。\par \textbf{详细解答} $2I=\iint[\sin x^2\cos y^2+\sin y^2\cos x^2]\dd A=\iint\sin(x^2+y^2)\dd A$。极坐标给$I=\pi\int_0^1r\sin(r^2)\dd r=\frac{\pi(1-\cos1)}{2}$。\par \textbf{验算与条件复核} 令$u=r^2$核对径向系数$\frac{1}{2}$。\par \textbf{疑难点与方法比较} 原式本身不是径向函数，不能不作轮换就直接把$x^2$写成$r^2$。\end{example}

\subsection{齐次边界的径向解析}\label{ic:delta_multiple:homogeneousboundary}\begin{method}{齐次边界的径向解析}\textbf{识别信号} $F(x,y)$是$d$次齐次函数，边界为$F=C$。\par \textbf{操作步骤} 代$x=r\cos\theta,y=r\sin\theta$，由$r^dF(\cos\theta,\sin\theta)=C$解径向上界；先确定右端为正的角区间，再积分。\par \textbf{疑难点与条件} 在角度分母为0或变号处必须分区；不可把负径向上界当成区域。\end{method}

\begin{example}{竞赛风格｜双叶区域的面积}{难度：$\star\star\star$}\label{ic:ex:delta_multiple:lemniscate}\textbf{题目} 求$(x^2+y^2)^2\le a^2(x^2-y^2)$的面积，$a>0$。\par \textbf{思路} 四次与二次齐次项使$r^2$可解。\par \textbf{详细解答} 极坐标给$0\le r^2\le a^2\cos2\theta$，只在$\cos2\theta\ge0$的角域存在非零点。右叶$\theta\in[-\frac{\pi}{4},\frac{\pi}{4}]$，面积$\frac{a^2}2\int_{-\frac{\pi}{4}}^{\frac{\pi}{4}}\cos2\theta\dd\theta=\frac{a^2}{2}$。左叶同面积，总面积$a^2$。\par \textbf{验算与条件复核} 极径上界在两端为0，且两叶只在原点接触，面积可相加。\par \textbf{疑难点与方法比较} 若误积分全角$[0,2\pi]$，正负抵消会错误得到0。\end{example}

\begin{example}{竞赛风格｜球内非径向矩的多坐标比较}{难度：$\star\star\star$}\label{ic:ex:delta_multiple:triplefourways}\textbf{题目} 求单位球内$\iiint(x^2+y^2)\dd V$。\par \textbf{思路} 可用轮换、柱坐标或截面。\par \textbf{详细解答} 轮换对称性给
\[\begin{aligned}
 I&=\dfrac23\iiint r^2\dd V\\
  &=\dfrac23\cdot\dfrac{4\pi}{5}=\dfrac{8\pi}{15}.
\end{aligned}\]
柱坐标独立计算为
\[\begin{aligned}
 I&=\int_0^{2\pi}\int_0^1\int_{-\sqrt{1-\rho^2}}^{\sqrt{1-\rho^2}}
       \rho^3\dd z\dd\rho\dd\theta\\
  &=4\pi\int_0^1\rho^3\sqrt{1-\rho^2}\dd\rho
   =\dfrac{8\pi}{15}.
\end{aligned}\]\par \textbf{验算与条件复核} 令$u=1-\rho^2$，末积分为$\frac12\int_0^1(1-u)u^{\frac{1}{2}}\dd u=\frac{2}{15}$。\par \textbf{疑难点与方法比较} 球坐标半径$r$和柱坐标半径$\rho$不同，不在同一式中混用。\end{example}

\section{图解：换序先读区域}
\begin{figure}[htbp]\centering
\begin{tikzpicture}[scale=4,>=stealth]
\fill[Teal!12](0,0) parabola(1,1)--(0,1)--cycle;
\draw[->](-.08,0)--(1.2,0)node[right]{$x$};
\draw[->](0,-.08)--(0,1.2)node[above]{$y$};
\draw[thick,Teal](0,0)parabola(1,1);
\draw[thick,Ink](0,1)--(1,1);
\draw[Gold,thick](.5,.25)--(.5,1);
\draw[Rust,thick](0,.5)--(.7071,.5);
\node[right]at(.8,.64){$y=x^2$};
\node[above]at(.55,1){$y=1$};
\node[below]at(1,0){$1$};
\end{tikzpicture}
\caption{同一区域$0\le x\le1,\ x^2\le y\le1$：竖切片给原积分限；横切片给$0\le y\le1,\ 0\le x\le\sqrt y$。颜色线段只展示一条典型切片，换序须覆盖所有切片。}\label{fig:review:order}
\end{figure}

\section{角向积分、联合分区与未知积分的代数化}\label{cmi:section}
\begin{examtip}
换元后仍要选择积分次序；分段表达式的分界线要共同决定区域划分。遇到未知函数被若干积分约束时，可以把这些积分视为待定常数，再建立有限维方程组。三重积分部分讨论积分区域的最优化，以及球面与抛物面相交时两种分层方式的核对。
\end{examtip}

\subsection{极坐标中先积角度}\label{cmi:angular}
\begin{method}{角向导数匹配优先于径向积分}
\textbf{识别信号} 换元后出现$\e^{r^2h(\theta)}h'(\theta)$、$\sin(r^2h(\theta))h'(\theta)$等结构。\par
\textbf{步骤} 先固定$r$对$\theta$积分，消去复合函数；再处理剩下的一元径向积分。\par
\textbf{疑难点} 极坐标并没有规定必须先对$r$积分，改变次序时仍须确认角度范围是否依赖$r$。
\end{method}
\begin{example}{竞赛风格｜先积角度消去复合指数}{难度：$\star\star\star$}
\label{cmi:ex:angular}
\textbf{题目} 在$D=\{x^2+y^2\le1,\ 0\le y\le x\}$上求
$I=\iint_D\e^{(x+y)^2}(x^2-y^2)\dd A$。\par
\textbf{思路} 极坐标使指数含$\sin2\theta$，而$x^2-y^2$提供其导数因子。\par
\textbf{详细解答} 令$x=r\cos\theta,y=r\sin\theta$，则$0\le r\le1$、$0\le\theta\le\frac{\pi}{4}$，
\[
I=\int_0^1r^3\dd r
\int_0^{\frac{\pi}{4}}\e^{r^2(1+\sin2\theta)}\cos2\theta\,\dd\theta.
\]
对$r>0$，内层积分等于
$[\frac{\e^{r^2(1+\sin2\theta)}}{2r^2}]_0^{\frac{\pi}{4}}$。所以
\[
\begin{aligned}
I&=\dfrac12\int_0^1r(\e^{2r^2}-\e^{r^2})\dd r\\
 &=\dfrac14\int_0^1(\e^{2u}-\e^u)\dd u
 =\dfrac{(\e-1)^2}{8},\qquad u=r^2.
\end{aligned}
\]
在$r=0$使用连续延拓理解中间表达式，单个端点也不影响积分。\par
\textbf{方法点评} 简化来自角向逆链式结构，而不是指数函数突然获得了初等径向原函数。\par
\textbf{易错点} 先把$\int\e^{r^2(1+\sin2\theta)}r^3\dd r$展开会增加运算量；必须保留原有$r^3$再与角积分系数合并。
\end{example}

\subsection{绝对值与最大值的共同分界}\label{cmi:partition}
\begin{method}{先叠加所有分界曲线，再逐块去符号}
\textbf{适用条件} 被积函数含绝对值、$\max$、$\min$或符号函数，且有多个分界条件。\par
\textbf{步骤} 列出每个表达式改变形式的曲线，判断它们在积分区间内的先后顺序；每块同时固定全部符号和分支。\par
\textbf{疑难点} 只按其中一条曲线划分，会使另一项在同一块内仍有两种表达式。
\end{method}
\begin{example}{竞赛风格｜带参数的三块积分}{难度：$\star\star\star$}
\label{cmi:ex:partition}
\textbf{题目} 对$p>1$，计算
\[
I(p)=\int_0^1\int_0^1|y-x^p|\max\{x,y\}\dd y\dd x.
\]
\textbf{思路} $0\le x\le1$时$x^p\le x$，因此$y=x^p$和$y=x$把正方形分为三块。\par
\textbf{详细解答} 依次在$0\le y\le x^p$、$x^p\le y\le x$、$x\le y\le1$积分：
\[
\begin{aligned}
I_1&=\int_0^1\int_0^{x^p}(x^p-y)x\dd y\dd x
    =\dfrac1{4(p+1)},\\
I_2&=\int_0^1\int_{x^p}^{x}(y-x^p)x\dd y\dd x\\
   &=\int_0^1\left(\dfrac{x^3}2-x^{p+2}+\dfrac{x^{2p+1}}2\right)\dd x\\
   &=\dfrac18-\dfrac1{p+3}+\dfrac1{4(p+1)},\\
I_3&=\int_0^1\int_x^1(y-x^p)y\dd y\dd x\\
   &=\int_0^1\left(\dfrac{1-x^3}3-\dfrac{x^p(1-x^2)}2\right)\dd x\\
   &=\dfrac14-\dfrac1{2(p+1)}+\dfrac1{2(p+3)}.
\end{aligned}
\]
相加时$(p+1)^{-1}$项完全抵消，故
\[
I(p)=\dfrac38-\dfrac1{2(p+3)}
=\dfrac{3p+5}{8(p+3)}.
\]
例如$p=2$时为$\frac{11}{40}$；令$p\downarrow1$，公式趋$\frac{1}{4}$，与直接按$y=x$两块积分相符。\par
\textbf{方法点评} 引入参数后，区域的先后关系和抵消结构比某一数值答案更清楚。\par
\textbf{易错点} 若$0<p<1$，两条曲线的顺序反转，不能原封不动使用这里的三个积分限。
\end{example}
\begin{figure}[htbp]\centering
\begin{tikzpicture}[x=5cm,y=4cm,>=stealth]
\fill[Teal!8](0,0)--(1,1)--(0,1)--cycle;
\fill[Gold!13]plot[domain=0:1,samples=65](\x,{\x*\x})--(1,1)--cycle;
\fill[Ink!6](0,0)--(1,0)--(1,1)
 --plot[domain=1:0,samples=65](\x,{\x*\x})--cycle;
\draw[Ink](0,0)rectangle(1,1);
\draw[Teal,thick](0,0)--(1,1);
\draw[Rust,thick]plot[domain=0:1,samples=65](\x,{\x*\x});
\draw[->,Ink](0,0)--(1.12,0)node[right]{$x$};
\draw[->,Ink](0,0)--(0,1.12)node[above]{$y$};
\node at(.22,.75){$D_3$};
\node at(.63,.51){$D_2$};
\node at(.73,.16){$D_1$};
\node[right]at(1,1){$y=x$};
\node[below]at(.48,.23){$y=x^2$};
\node[below left]at(0,0){$O$};
\end{tikzpicture}
\caption{$p=2$时的共同分区。跨过抛物线改变绝对值符号，跨过对角线改变最大值的取值；两项必须同时处理。}
\label{cmi:fig:partition}
\end{figure}

\subsection{有限个未知积分化为代数方程}\label{cmi:finiterank}
\begin{method}{先把积分记成常数，再积分恢复常数}
\textbf{适用条件} 未知函数通过有限个定积分或重积分进入等式，而且这些积分对自由变量为常数。\par
\textbf{步骤} 分别记$A,B,\ldots$，将未知函数写成显式函数与这些常数的线性组合；沿原有积分算子再次积分，解有限维线性方程组，最后代回验算。\par
\textbf{疑难点} 方程组可能无解或有多个解；不能默认每个积分关系都会唯一确定函数。
\end{method}
\begin{example}{考研风格｜含一维迹积分和二重积分的函数方程}{难度：$\star\star\star$}
\label{cmi:ex:finiterank}
\textbf{题目} $f$在$D=[0,1]^2$连续，并满足
\[
f(x,y)=xy+x+2\int_0^1f(t,0)\dd t+\iint_Df(u,v)\dd u\dd v.
\]
求$f$。\par
\textbf{思路} 两个积分与$x,y$无关；先合并成一个未知常数。\par
\textbf{详细解答} 记
$A=\int_0^1f(t,0)\dd t$、$B=\iint_D f$、$C=2A+B$，则$f=xy+x+C$。
沿$y=0$积分以及在单位正方形上积分，分别给
\[
A=\dfrac12+C,\qquad B=\dfrac34+C.
\]
所以
$C=2(\frac12+C)+(\frac34+C)=\frac74+3C$，
解得$C=-\frac{7}{8}$。因此
\[
f(x,y)=xy+x-\dfrac78,\qquad A=-\dfrac38,\quad B=-\dfrac18.
\]
代回$2A+B=-\frac{7}{8}$，原关系成立；而$C$的线性方程有唯一解，故$f$也唯一。\par
\textbf{方法点评} 二重积分作为一个数参与方程，避免把问题误当成偏微分方程。\par
\textbf{易错点} 若把系数改变到使常数方程的左右$C$系数相同，需要另查相容性，不能直接相除。
\end{example}

\subsection{按正负选择积分域与球抛面交域}\label{cmi:regions}
\begin{theorem}{不受体积约束时的正部选域}
若实函数$h$的正部$h_+=\max(h,0)$在全空间可积，则对任何使积分有意义的可测区域$\Omega$，
\[
\int_\Omega h\le\int_{\R^d}h_+.
\]
当允许取$\{h>0\}$时，上界可以达到。等号要求未选入的正值部分与选入的负值部分的积分均为零。
\end{theorem}
\begin{proof}
设$E=\{h>0\}$，则
\[
\int_Eh-\int_\Omega h
=\int_{E\setminus\Omega}h+\int_{\Omega\setminus E}(-h)\ge0.
\]
两项均非负，因而也给出等号条件。
\end{proof}
\begin{method}{先判断是否有选域约束}
\textbf{适用条件} 题目要求在区域变化时使积分最大，且没有固定体积、固定边界等额外限制。\par
\textbf{步骤} 选择全部正值区，排除负值区；再用线性缩放或合适坐标计算最大值。\par
\textbf{疑难点} 若限制体积固定，通常要选某个超水平集，不能直接把整个正值区全部选入。
\end{method}
\begin{example}{竞赛风格｜同时求最优三维区域与最大值}{难度：$\star\star\star$}
\label{cmi:ex:positivepart}
\textbf{题目} 在有界分段光滑区域$\Omega\subset\R^3$中选择区域，使
$I(\Omega)=\iiint_\Omega(1-x^2-2y^2-3z^2)\dd V$最大，并求该最大值。\par
\textbf{思路} 先用正部选域，不对未知边界作变分。\par
\textbf{详细解答} 可取
$\Omega_*=\{x^2+2y^2+3z^2<1\}$。
作$x=u,y=\frac{v}{\sqrt2},z=\frac{w}{\sqrt3}$，Jacobian绝对值$\frac{1}{\sqrt6}$，得到
\[
\begin{aligned}
I_{\max}
&=\dfrac1{\sqrt6}\iiint_{u^2+v^2+w^2<1}(1-u^2-v^2-w^2)\dd u\dd v\dd w\\
&=\dfrac{4\pi}{\sqrt6}\int_0^1(1-r^2)r^2\dd r
=\dfrac{8\pi}{15\sqrt6}.
\end{aligned}
\]
椭球边界上被积式为零，取开域或加入边界不改变积分。\par
\textbf{方法点评} 选域依据是符号，坐标换元只用于计算已确定的最大值。\par
\textbf{易错点} 这里没有预先固定体积；如果增加体积条件，最优区域的推导须重新进行。
\end{example}

\begin{example}{考研风格｜球面与抛物面之间两种分层}{难度：$\star\star\star$}
\label{cmi:ex:sphereparaboloid}
\textbf{题目} $\Omega$由抛物面$z=x^2+y^2$和上半球面$z=\sqrt{2-x^2-y^2}$围成，计算$\iiint_\Omega z\dd V$。\par
\textbf{思路} 柱坐标适合先对$z$积分；水平切片则在交线高度处分段。\par
\textbf{详细解答} 令$r^2=x^2+y^2$。交线满足$r^2=\sqrt{2-r^2}$，
即$(r^2-1)(r^2+2)=0$，所以$r=1,z=1$。
柱坐标给
\[
\begin{aligned}
I&=2\pi\int_0^1r\dd r\int_{r^2}^{\sqrt{2-r^2}}z\dd z\\
 &=\pi\int_0^1(2r-r^3-r^5)\dd r=\dfrac{7\pi}{12}.
\end{aligned}
\]
第二法固定$z$：在$0\le z\le1$时截面为$r^2\le z$，面积$\pi z$；
在$1\le z\le\sqrt2$时截面为$r^2\le2-z^2$，面积$\pi(2-z^2)$。
因此
\[
I=\pi\int_0^1z^2\dd z+\pi\int_1^{\sqrt2}z(2-z^2)\dd z
=\dfrac{\pi}3+\dfrac{\pi}4=\dfrac{7\pi}{12}.
\]
\textbf{方法点评} 同一体域的两个截面描述核对了球面根号、交线和积分限。\par
\textbf{易错点} 上边界是$\sqrt{2-r^2}$，不是$2-r^2$；水平截面必须在$z=1$处改变控制边界。
\end{example}

\section*{章末分层习题}\addcontentsline{toc}{section}{章末分层习题}
\refstepcounter{exercise}\label{q:7:1}\par\noindent\textbf{习题\theexercise\quad 基础题} 求单位正方形上$\iint xy\dd A$。\quad\hyperref[a:7:1]{[答案]}\par
\refstepcounter{exercise}\label{q:7:2}\par\noindent\textbf{习题\theexercise\quad 基础题} 求半径2圆盘面积。\quad\hyperref[a:7:2]{[答案]}\par
\refstepcounter{exercise}\label{q:7:3}\par\noindent\textbf{习题\theexercise\quad 基础题} 写出球坐标的体积元。\quad\hyperref[a:7:3]{[答案]}\par
\refstepcounter{exercise}\label{q:7:4}\par\noindent\textbf{习题\theexercise\quad 考研风格} 求单位圆盘上$\iint x\dd A$。\quad\hyperref[a:7:4]{[答案]}\par
\refstepcounter{exercise}\label{q:7:5}\par\noindent\textbf{习题\theexercise\quad 考研风格} 求单位球内$\iiint(x^2+y^2+z^2)\dd V$。\quad\hyperref[a:7:5]{[答案]}\par
\refstepcounter{exercise}\label{q:7:6}\par\noindent\textbf{习题\theexercise\quad 考研风格} 设$f$在闭区域$D=\{0\le x\le1,\ x^2\le y\le1\}$上连续，将$\int_0^1\dd x\int_{x^2}^1 f(x,y)\dd y$换序。\quad\hyperref[a:7:6]{[答案]}\par
\refstepcounter{exercise}\label{q:7:7}\par\noindent\textbf{习题\theexercise\quad 竞赛拓展} 求$\iint_{\R^2}\e^{-2(x^2+y^2)}\dd A$。\quad\hyperref[a:7:7]{[答案]}\par
\refstepcounter{exercise}\label{q:7:8}\par\noindent\textbf{习题\theexercise\quad 竞赛拓展} 设$a,b,c>0$，求椭球$\frac{x^2}{a^2}+\frac{y^2}{b^2}+\frac{z^2}{c^2}\le1$的体积。\quad\hyperref[a:7:8]{[答案]}\par
\refstepcounter{exercise}\label{q:7:9}\par\noindent\textbf{习题\theexercise\quad 竞赛拓展} 设$f$为闭单位圆盘$D$上的实连续函数，证明$\iint_D f^2\dd A\ge\frac{(\iint_D f\dd A)^2}{\pi}$。\quad\hyperref[a:7:9]{[答案]}\par
\section*{章末方法速查}\addcontentsline{toc}{section}{章末方法速查}
\par\noindent\begin{tabularx}{\textwidth}{p{23mm}XXX}\toprule 方法&适用条件&典型信号&易错点\\\midrule
换序&区域可重述&内积分困难&边界漏分块\\极坐标&圆与径向&$x^2+y^2$&漏$r$\\球坐标&球锥区域&半径与角&漏$\sin\phi$\\一般换元&一一对应&线性边界&Jacobian倒置\\
\bottomrule\end{tabularx}
